\subsection{Why a literal witness reuse does not shorten this emitter}
A natural shortening attempt is to reuse an already emitted head as the
required arbitrary diagonal clause, thereby removing the eight-token
construction instead of moving it.  There is a small syntactic obstruction
for the current axiom families.

The numerical implication constructor is injective.  Its output is
congruent to zero modulo four, whereas a universal formula has tag two.
Moreover $U(x,Y)>\max(x,Y)$ for every natural pair, including zero.
Thus $\Gamma(B)=I(B,U(B,B))$ cannot coincide with heads 3, 4, or 6, or
with the dominated propositional head: each has an implication, of tag zero,
as its outer consequent.  Head 5 and the unbundled gate have outer consequent
$1$, which likewise differs from $U(B,B)\ge2$.

Only heads 1 and 2 require an additional check.  If
\[
 I(\exists x_d U(d,h),h)=I(B,U(B,B)),
\]
then $h=U(B,B)>B$, while
$B=\exists x_d U(d,h)>h$, a contradiction.  If
\[
 I(U(d,U(a,h)),U(a,U(d,h)))=I(B,U(B,B)),
\]
then universal-constructor injectivity gives $a=B$ and hence
$B=U(d,U(B,h))>B$, again a contradiction.  All inequalities are elementary
PA arithmetic facts about the displayed Cantor codes.  The zero-code alias
of implication does not evade them, since the universal terms are positive
and strictly increasing in the sense just used.

This excludes literal coincidence with the retained heads, not derivability
from them and not a smaller machine by another construction.  If the
parameter $B$ is restricted to formulas already known to be theorems, then
$\Gamma(B)$ is already obtained by closed generalization followed by
implication introduction.  Such a reuse supplies no uniform replacement,
by itself, for the arbitrary negative-formula instances used in the PA195
reverse-equality proof.  A genuine shortening through this route therefore
requires a new calibration argument or a different arithmetic emitter.
